% English counterpart of the focal Spanish insertion supplied by Ley 9 puertas.
% Provenance and verification are retained in the associated handoff package.
\subsubsection{Calendar recovery and ternary carry}
\label{ins29:xii:acarreo}

The extraction of the quarter-turn plane admits an exact description
of the temporal information retained. We use the calendar of
\(108\) positions and the dodecaphasic extension over the nonadic phase
constructed earlier. For this description we use zero-based representatives:
\[
 \theta=r+9\beta,\qquad 0\leq r<9,\quad 0\leq\beta<12.
\]
The position enumerated from \(1\) to \(9\) in the record is \(r+1\) here.
The enriched update precedes this finite calendar and also retains
the route, incidences, and accumulated memory. On the plane already
constructed, \(C_{12}^{\beta}W=WJ_0^{\beta}\); its reading therefore
depends on \(b=\beta\bmod4\).

\begin{proposition}[Ternary fibre of the quarter-turn reading]
\label{ins29:xii:fibra}
The pair consisting of the nonadic phase \(r\) and the planar phase \(b\)
determines a quotient of \(36\) calendar positions. Its composition law is
\[
 (r,b)(r',b')=
 \left(r+r'\bmod9,\;
 b+b'+\left\lfloor\frac{r+r'}9\right\rfloor\bmod4\right).
\]
Each reading has three preimages in \(\mathbb Z/108\mathbb Z\).
The additional component \(k=\beta\bmod3\) recovers the dodecaphasic
coordinate through
\[
 \beta=9b+4k\pmod{12}.
\]
\end{proposition}
\begin{proof}
Adding two representatives gives
\(\theta+\theta'=r+r'+9(\beta+\beta')\). Euclidean division of
\(r+r'\) by \(9\) produces the carry in the formula. The reading
\((r,b)\) is equivalent to the residue \(\theta\bmod36\), since
\(\theta\equiv r+9b\pmod{36}\). Its kernel is
\(\{0,36,72\}\), so each fibre has three elements.
Finally, \(9b+4k\) is congruent to \(b\) modulo \(4\) and to
\(k\) modulo \(3\); these two residues determine a unique element
modulo \(12\).
\end{proof}

\begin{proposition}[Localization of the non-split extension]
\label{ins29:xii:extension}
The decomposition
\[
 \theta\longmapsto(j,z)=(\theta\bmod4,\theta\bmod27),
 \qquad \theta=81j+28z\pmod{108},
\]
identifies the calendar with
\(\mathbb Z/4\mathbb Z\times\mathbb Z/27\mathbb Z\).
The nonadic projection is \((j,z)\mapsto z\bmod9\), and the
dodecaphasic inclusion \(\beta\mapsto9\beta\) takes the form
\[
 \beta\longmapsto(\beta\bmod4,9\beta\bmod27).
\]
The obstruction to splitting the full extension is concentrated in
\[
 0\longrightarrow\mathbb Z/3\mathbb Z
 \longrightarrow\mathbb Z/27\mathbb Z
 \longrightarrow\mathbb Z/9\mathbb Z\longrightarrow0.
\]
\end{proposition}
\begin{proof}
The coefficients \(81\) and \(28\) have residues \((1,0)\) and \((0,1)\)
modulo \((4,27)\), respectively. This verifies the inverse and the two
specified maps. The ternary sequence does not split: a product
\(\mathbb Z/3\mathbb Z\times\mathbb Z/9\mathbb Z\) has exponent
\(9\), whereas \(\mathbb Z/27\mathbb Z\) has exponent \(27\).
By contrast, the projection of the \(36\)-position quotient onto
\(\mathbb Z/9\mathbb Z\) admits the homomorphic section
\(n\mapsto28n\pmod{36}\). Thus, the reduction removes precisely
this ternary obstruction.
\end{proof}

The distinction is already visible at \(\theta=0,36,72\): all three
positions give \(r=b=0\), whereas \(z\) takes the values \(0,9,18\),
respectively. Recovering the calendar requires retaining the fibre
identified by the plane. Moreover, the global quaternary residue satisfies
\(j=r+b\pmod4\); the planar phase \(b\) is transported together with
the nonadic phase and its carry.

This locates the scope of the quarter turn: it organizes an oriented
representation of the record and, together with the nonadic residue,
recovers exactly its quotient of \(36\) positions. The additional ternary
coordinate restores the \(108\) positions. The memory of the enriched
dynamics in turn retains the continuation between complete calendar returns.
The periodicity of a representation and the persistence of the recorded
trajectory therefore belong to distinct readings of the same transport.
